% =====================================================================
%  Thème 4 — CORRIGÉ — Niveau MOYEN
% =====================================================================
\documentclass[11pt,a4paper]{article}
\usepackage[utf8]{inputenc}
\usepackage[T1]{fontenc}
\usepackage{lmodern}
\usepackage[french]{babel}
\usepackage{amsmath,amssymb,mathtools}
\usepackage{icomma}
\usepackage{xcolor}
\usepackage[most]{tcolorbox}
\usepackage{enumitem}
\usepackage{array}
\usepackage{booktabs}
\usepackage{fancyhdr}
\usepackage[hidelinks]{hyperref}
\usepackage[a4paper,margin=2.2cm,headheight=26pt,headsep=10pt]{geometry}

\definecolor{primary}{RGB}{223,87,99}
\definecolor{primarydark}{RGB}{150,42,55}
\definecolor{excol}{RGB}{0,135,90}

\tcbset{boxbase/.style={enhanced,breakable,boxrule=0.9pt,arc=2.5pt,
   left=3mm,right=3mm,top=2.3mm,bottom=2.3mm,
   fonttitle=\bfseries,coltitle=white,toptitle=0.6mm,bottomtitle=0.6mm}}
\newtcolorbox[auto counter]{correction}[1][]{boxbase,colback=excol!5,colframe=excol,
  title={\textsc{Corrigé}~\thetcbcounter\ifx\relax#1\relax\relax\else\textmd{ --- \itshape #1}\fi}}

\pagestyle{fancy}
\fancyhf{}
\fancyhead[L]{\small\textcolor{primarydark}{\textbf{Fiche 6} — Les limites}}
\fancyhead[R]{\small\textcolor{primarydark}{\textsc{Thème 4 — Corrigé Moyen}}}
\fancyfoot[C]{\small\thepage}
\renewcommand{\headrulewidth}{0.8pt}
\renewcommand{\headrule}{\hbox to\headwidth{\color{primary}\leaders\hrule height \headrulewidth\hfill}}

\newcommand{\R}{\mathbb{R}}

\begin{document}

\begin{center}
{\Large\bfseries\color{primarydark}Thème 4 — Radicaux : facteur dominant \& conjugué}\\[2pt]
{\large\color{excol}Corrigé — Niveau Moyen}
\end{center}
\vspace{6pt}

\begin{correction}[le signe de $|x|$ change tout]
$f(x)=\dfrac{\sqrt{x^2+1}}{x}$.
\begin{enumerate}[label=\alph*),leftmargin=2em,itemsep=3pt]
  \item $\sqrt{x^2+1}=\sqrt{x^2\bigl(1+\tfrac{1}{x^2}\bigr)}=|x|\sqrt{1+\tfrac{1}{x^2}}$.
  \item En $+\infty$, $|x|=x$ : $\displaystyle f(x)=\frac{x\sqrt{1+\frac{1}{x^2}}}{x}=\sqrt{1+\tfrac{1}{x^2}}\to\sqrt1=1$.
        Donc $\displaystyle\lim_{x\to+\infty} f(x)=1$.
  \item En $-\infty$, $|x|=-x$ : $\displaystyle f(x)=\frac{-x\sqrt{1+\frac{1}{x^2}}}{x}=-\sqrt{1+\tfrac{1}{x^2}}\to-1$.
        Donc $\displaystyle\lim_{x\to-\infty} f(x)=-1$. \emph{(Deux limites différentes : oublier le $|x|$ aurait donné $1$ à tort.)}
\end{enumerate}
\end{correction}

\begin{correction}[différence de racines : $\infty-\infty$]
\begin{enumerate}[label=\alph*),leftmargin=2em,itemsep=3pt]
  \item $\sqrt{x^2+x}\to+\infty$ et $x\to+\infty$ : forme $\infty-\infty$ (indéterminée).
  \item $\displaystyle\sqrt{x^2+x}-x=\frac{(x^2+x)-x^2}{\sqrt{x^2+x}+x}=\frac{x}{\sqrt{x^2+x}+x}$.
  \item Pour $x>0$, $\sqrt{x^2+x}=x\sqrt{1+\tfrac1x}$, donc
        \[ \frac{x}{x\sqrt{1+\frac1x}+x}=\frac{x}{x\bigl(\sqrt{1+\frac1x}+1\bigr)}=\frac{1}{\sqrt{1+\frac1x}+1}
        \xrightarrow[x\to+\infty]{}\frac{1}{1+1}=\frac12. \]
\end{enumerate}
\end{correction}

\begin{correction}[$\frac00$ en un point]
En $3$ : $\sqrt{4}-2=0$ et $3-3=0$ ($\frac00$). On multiplie par $\sqrt{x+1}+2$ :
\[ \frac{\sqrt{x+1}-2}{x-3}=\frac{(x+1)-4}{(x-3)\bigl(\sqrt{x+1}+2\bigr)}
  =\frac{x-3}{(x-3)\bigl(\sqrt{x+1}+2\bigr)}=\frac{1}{\sqrt{x+1}+2}. \]
Donc $\displaystyle\lim_{x\to 3}\frac{\sqrt{x+1}-2}{x-3}=\frac{1}{\sqrt{4}+2}=\frac{1}{4}$.
\end{correction}

\begin{correction}[somme de racine et polynôme]
\begin{enumerate}[label=\alph*),leftmargin=2em,itemsep=3pt]
  \item Quand $x\to-\infty$ : $\sqrt{x^2+1}\to+\infty$ mais $x\to-\infty$ ; la somme est du type
        $(+\infty)+(-\infty)$, soit $\infty-\infty$ (indéterminée).
  \item $\displaystyle\sqrt{x^2+1}+x=\frac{(x^2+1)-x^2}{\sqrt{x^2+1}-x}=\frac{1}{\sqrt{x^2+1}-x}$.
        Or, en $-\infty$, $\sqrt{x^2+1}\to+\infty$ et $-x\to+\infty$, donc le dénominateur $\to+\infty$ :
        \[ \lim_{x\to-\infty}\bigl(\sqrt{x^2+1}+x\bigr)=\frac{1}{+\infty}=0. \]
\end{enumerate}
\end{correction}

\begin{correction}[racine sur quotient factorisable]
\begin{enumerate}[label=\alph*),leftmargin=2em,itemsep=3pt]
  \item $x^2-4=(x-2)(x+2)$.
  \item En multipliant par le conjugué $\sqrt{x+2}+2$ :
  \[ \frac{\sqrt{x+2}-2}{x^2-4}=\frac{(x+2)-4}{(x-2)(x+2)\bigl(\sqrt{x+2}+2\bigr)}
     =\frac{x-2}{(x-2)(x+2)\bigl(\sqrt{x+2}+2\bigr)}=\frac{1}{(x+2)\bigl(\sqrt{x+2}+2\bigr)}. \]
  En $x=2$ : $\displaystyle\frac{1}{4\cdot(\sqrt4+2)}=\frac{1}{4\cdot 4}=\frac{1}{16}$.
\end{enumerate}
\end{correction}

\end{document}
